\documentclass[11pt,a4paper]{article}
\usepackage[T1]{fontenc}
\usepackage[utf8]{inputenc}
\usepackage[polish]{babel}
\usepackage{lmodern}
\usepackage[a4paper,margin=2.3cm]{geometry}
\usepackage{amsmath,amssymb,booktabs}
\title{HF20: matematyka horyzontów i bramy jakości opadu}
\author{Edward Szczypka}
\date{\today}
\begin{document}
\maketitle
\section{Horyzont użytkownika i modelu}
\[
 d=\left\lceil\frac{a-t_o}{1\,\mathrm{h}}\right\rceil,
 \qquad h_\ell=d+\ell-1,
 \qquad \ell=1,\ldots,48.
\]
Dla $d\leq 13$ otrzymujemy $h_{48}\leq 60$.
\section{Model hurdle opadu}
Niech $Z=\mathbf{1}\{Y\geq\tau\}$. Model szacuje
$p(x)=\Pr(Z=1\mid x)$ oraz $m(x)=\mathbb{E}[Y\mid Z=1,x]$, a prognoza
oczekiwana wynosi $\widehat Y=p(x)m(x)$.
\section{Brier skill}
Dla $BS=\frac1n\sum_i(p_i-z_i)^2$ i baseline $BS_0$:
\[
 BSS=1-\frac{BS}{BS_0}.
\]
Wymagamy dodatniego BSS względem klimatologii i persistence, poprawy MAE
względem persistence, ograniczonego biasu oraz ROC AUC nie mniejszego niż
ustalony próg.
\end{document}
